% claim-ref: guppy-grass#method
% claim-ref: guppy-grass#contain-discard
% claim-ref: guppy-grass#observe-success
\CompanionHeader{PROPAGATION}{Guppy grass}{Najas guadalupensis}{Aquatic plant / submerged stem fragments. Proposed home sequence; identity and tank readings pending.}
\Context{Start only with healthy material and a prepared destination. Keep aquatic tissue wet; do not import terrestrial soil, callusing or watering schedules.}
\begin{minipage}[t]{.625\linewidth}\raggedright\vspace{0pt}
\Step{1}{Prepare / contain}{Use clean scissors, a container of tank water and a waste bag. Inspect for unwanted material; keep clippings contained. \textbf{[DEC]}}
\Step{2}{Select}{Choose an intact growing stem from the reported guppy grass. Verify the working identity before treating this as species-specific guidance. \textbf{[GG-TAM]}}
\Step{3}{Separate}{Gently separate a leafy fragment. Pond evidence establishes regrowth capability, not a universal cutting length or time to roots. \textbf{[GG-PROP]}}
\Step{4}{Place floating}{Keep the cutting submerged as a loose floating mass. No glue, terrestrial callus period or rooting hormone is needed for this proposed floating method. \textbf{[GG-AA]}}
\Step{5}{Maintain / observe}{Keep recorded livestock-compatible conditions. Look for new shoots and sustained intact growth; one detached fragment is not proof of establishment. \textbf{[GG-AA, GG-PROP]}}
\Step{6}{Thin / record}{Remove crowding and collect loose stems. Leave equipment and animal access clear; log growth and conditions before changing them. \textbf{[GG-MSU]}}
\Step{7}{Dispose responsibly}{Seal unwanted fragments for trash; never release or compost aquarium plants. \textbf{[DEC]}}
\end{minipage}\hfill\begin{minipage}[t]{.345\linewidth}\raggedright\vspace{0pt}
\Note{Milestones}{Separate material; observe new growth; confirm continued healthy growth. A cut or glue bond is only the starting event. No universal days-to-success is established.}
\Note{If growth stalls}{Record light, water, flow and recent changes. Inspect for detached, buried or deteriorating tissue. These are checks, not a diagnosis or medication instruction.}
\Note{No safety guarantee}{Plants do not prove cycling, sufficient oxygen or predation-safe shelter. Review animal care and current water results separately.}
\end{minipage}\par\vspace{4pt}
\Panel{Observation record}{Start / method: \hrulefill\quad Date: \hrulefill\par New growth / attachment / conditions: \hrulefill}
\Sources{\href{https://aquaticarts.com/products/guppy-grass}{GG-AA: Aquatic Arts}; \href{https://aquaplant.tamu.edu/plant-identification/alphabetical-index/southern-naiad/}{GG-TAM: Texas A\&M AgriLife Extension / AquaPlant}; \href{https://aquaplant.tamu.edu/management-options/southern-naiad/}{GG-PROP: Texas A\&M AgriLife Extension / AquaPlant}; \href{https://extension.msstate.edu/publications/southern-naiad-najas-guadalupensis}{GG-MSU: Mississippi State University Extension}; \href{https://www.calflora.org/app/taxon?crn=5743}{GG-CAL: Calflora}; \href{https://dec.ny.gov/nature/animals-fish-plants/invasive-species/aquatic/prevent-spread-of-aquatic-invasive-species/guidelines-for-aquarium-and-pet-owners}{DEC: New York State Department of Environmental Conservation}. Full scopes and claim records: adjacent YAML. Proposed instructions adapt cited methods; local conditions remain unknown.}{guppy-grass / propagation}
